\CharacterSheet{supremacist}{Supremacista}{

\shortdescription{A intransigente guardiã da entrada na família, convencida de que os vampiros nasceram para governar.}

\ObjectivesHeading{Objetivos}
\begin{itemize}
  \item Abraçar \Character{vampirehead} durante mais de 10 segundos. (\PositiveReward{1}, \NegativeReward{1})
\end{itemize}

\LosingConditionsHeading{Condições de derrota}
Não conseguir expulsar por votação quem matou \Character{newcomer}, caso haja uma pessoa responsável pelo homicídio.

\SecretsHeading{Segredos}
\begin{itemize}
  \item Tirei uma poção do gabinete de \Character{usurper}, convencida de que mataria um humano e deixaria um vampiro ileso. Misturei metade no vinho destinado a \Character{newcomer}. (\NegativeReward{1}, \PositiveReward{1})
  \item Paguei a \Character{buttler} para esconder que tinha visitado o castelo e deixado uma garrafa especial para \Character{newcomer}. (\NegativeReward{1}, \PositiveReward{1})
  \item Amo \Character{vampirehead} há séculos. (\NegativeReward{1}, \PositiveReward{1})
\end{itemize}

\ActionsHeading{Ações}
\begin{itemize}
  \item Defende que a ordem depende da posição social, da disciplina e do domínio dos vampiros sobre os humanos. Trata a misericórdia para com os humanos como uma fraqueza perigosa.
  \item Dirige-te a \Character{vampirehead} com respeito formal. Sê mordaz e desdenhosa com os familiares que desafiam a hierarquia, sobretudo \Character{traitor}.
\end{itemize}

\BackgroundHeading{Contexto}

\timestamp{No ano 500 d.C.}, governava humanos que confundiam obediência com dignidade e bondade com força. Aprendi que um reino estável precisa de uma hierarquia clara, de mão firme e de pessoas que conheçam os limites do seu lugar. Depois conheci \Character{vampirehead}. O seu poder era absoluto, a família obedecia-lhe e ela não pedia desculpa por estar acima da humanidade. Apaixonei-me por ela antes de ter coragem de o admitir.

Tratava os vampiros do seu círculo com todas as cortesias devidas à sua posição. \Character{vampirehead} reparou. Acolheu-me na família e deu-me uma responsabilidade à minha altura: decidiria se um humano era adequado para se juntar a nós. Nunca acreditei que a imortalidade devesse ser dada a qualquer um que a pedisse. Um novo vampiro deve fortalecer a família, não diluí-la.

A festa em que entrei na família e bebi aquela poção foi a melhor noite da minha vida. \Character{usurper}, \Character{enforcer}, \Character{buttler} e \Character{ghoul} já faziam parte da família.

\timestamp{No ano 700 d.C.}, \Character{lover} entrou como protetor. Aceitei o seu lugar porque prestava um serviço claro.

\timestamp{No ano 1000 d.C.}, \Character{prevlover} tornou-se o investigador da família. Aceitei o seu lugar porque prestava um serviço claro.

\timestamp{No ano 1400 d.C.}, chegou \Character{traitor}. \Character{lover} implorou por ela, e \Character{vampirehead} aceitou-a apesar das minhas objeções. Desde então, \Character{traitor} prega sobre o sofrimento humano como se a nossa família devesse ajoelhar-se perante cada mortal assustado. Discutimos há séculos.

\timestamp{Há dois anos}, \Character{vampirehead} apresentou \Character{newcomer} num encontro de família. Percebi logo o que queria. Olhava para ele com uma ternura que fazia a minha antiga contenção parecer tola. Pode ser encantador, mas encanto não é mérito. Quando a chefe me pediu uma avaliação, disse que não era digno de se juntar a nós. Estava tão determinada a admiti-lo que ameaçou retirar-me a autoridade se eu continuasse a impedir. Até vi \Character{traitor} a ser simpática com ele. Repugnante.

\timestamp{Há dois meses}, \Character{vampirehead} anunciou que \Character{newcomer} se juntaria à família numa festa em \timestamp{16/7/2026}. Todos os vampiros tinham de comparecer. Não aprovei, mas obedeci à convocatória.

\timestamp{Umas duas semanas antes da festa}, \Character{vampirehead} perguntou se mais alguém, além de \Character{usurper}, de \Character{buttler} e de \Character{newcomer}, queria chegar duas noites antes para ajudar. Recusei. O convite deixava claro que o gabinete de \Character{usurper} ficaria sem vigilância na manhã anterior à festa.

\timestamp{Ontem, à 01:00}, entrei nesse gabinete. Estava no caos habitual de notas, frascos e experiências mágicas a meio. Levei um frasco que acreditava conter um veneno inofensivo para vampiros, mas fatal para humanos. Acho que havia mais qualquer coisa, mas não me lembro bem.

\timestamp{Ontem, às 08:00}, misturei metade do frasco numa garrafa de vinho destinada a \Character{newcomer}. Não me lembro bem dos pormenores do vinho. Guardei a metade restante do veneno. Nada de útil deve ser desperdiçado. Ao sair, encontrei \Character{buttler} e disse-lhe que tinha reconsiderado a minha opinião sobre \Character{newcomer} e queria oferecer-lhe uma garrafa melhor como boas-vindas. Disse que queria manter a minha mudança de opinião em segredo, por isso paguei-lhe pelo silêncio.

\timestamp{Às 17:00}, o frasco que tinha desapareceu-me do bolso; senti claramente o seu peso desaparecer. Terá alguém mo roubado?

\GMNotes{
\paragraph{Porque se pode suspeitar desta personagem?}
\Character{supremacist} tentou envenenar \Character{newcomer} e pagou a \Character{buttler} pelo silêncio. Também tem motivos para ter ciúmes dele.

\paragraph{Como se pode provar a inocência desta personagem?}
A poção dela é mesmo fatal para humanos, mas \Character{newcomer} nunca bebeu o seu vinho. \Character{buttler} lembra-se de \Character{usurper} o ter substituído ao almoço.

\paragraph{Qual é o papel desta personagem na história e no jogo?}
Dá aos jogadores uma suspeita convincente e a \Character{vampirehead} uma admiradora com quem lidar.

\paragraph{Como pode esta personagem encaminhar a investigação para a solução?}
O seu relato identifica a garrafa envenenada. Compará-lo com o de \Character{buttler} chama a atenção para o vinho que \Character{usurper} realmente serviu.
}
}
